\section{Discussion}
\label{sec:discussion}

The strict positives show that a connectivity relation can, in a restricted
event class, outlast complete replacement of the cells that instantiate its
bridge. The particular bridge is therefore not sufficient to identify what
persists in these events. The invariant candidate is the restored reachability
relation between tracked region descendants.

This observation is consistent with the organizing hypothesis: interaction can
alter the set of reachable collective states, and a constraint on that set can
be re-instantiated in a different local realization. The result does not by
itself validate the general reachable-set information measure in
Equation~\ref{eq:information}.

Findings 066--067 add an event-resolved mechanism. A prospective expression can
survive the death of its stationary address when an ancestry-preserving
successor is realized while that route remains available. A lawful route may be
lost because a gate fails, the source spends its final amplitude elsewhere, or
another source activates the target first. The operational content of
\textit{before} is therefore causal rather than merely chronological: changing
one identified decision changes whether the route remains available or whether
a child is realized.

The intervention results also prevent a stronger but unsupported conclusion.
Route availability and realization are distinct. Every exhaustion intervention
in the prospective same-state cohort reopened the route, yet two did not later
hand off; the independent seed-58 exhaustion case likewise reopened opportunity
without rescue. Information, in the organizing language, concerns a structured
set of still-possible futures whose members have lawful lifetimes. Persistence
requires both an open route and timely realization. No payload transport is
needed for the operational lineage used here.

Finding 068 shows that the naturally realized child is not generally unique in
the reachable structure. Although every factual branch expressed exactly one
child, veto exposed an alternate realization in 122 of 164 counterfactuals,
usually at a different target and usually after a delay. The particular child
was therefore temporally redundant in those branches. In the remaining 42, no
alternate child appeared and the lineage died at address loss, establishing
branch-level but-for necessity for the selected realization.

Finding 069 shows that this redundancy can recur through several ordered
replacements. Observed policy cuts occurred at every tested depth from one
through eight: \GCTenCutTotal{} of \GCTenInitial{} branches lost the lineage
after some committed prefix, whereas \GCTenAboveMax{} still handed off after
all eight vetoes and \GCTenSourceCensored{} retained the original address at the
fixed horizon. The depth is therefore a measure of robustness to this declared
sequence of realization suppressions. It is not evidence that a static reserve
contained that many interchangeable children. Each veto retained amplitude and
changed the subsequent lawful path, on which a later successor could become
available and be realized. Target reuse in \GCTenTargetReuseBranches{} branches
also reinforces the distinction between a stationary address and the lineage
relation expressed through it.

Together these results support a sharper hierarchy. Information describes
which futures remain possible; constraint determines which are lawful; ordered
interaction gives those routes finite opportunity lifetimes; realization
selects one route; and redundancy determines whether another realization can
substitute when that selection is removed. The natural trajectory alone
concealed most of this redundancy because it displayed one child even when a
different future could become reachable along the retained-amplitude
counterfactual path.

Finding 070 moves that conclusion beyond the original selected source states.
Every one of four prospectively screened independent geometries admitted an
alternate continuation after the first veto, and cuts again occurred at every
tested depth. This supports geometric recurrence of ordered continuation
redundancy under the frozen law and scheduler. It does not turn the adaptive
chains into a static map: target reuse and retained amplitude mean that later
routes can revisit addresses or arise only in the altered history.

Finding 074 directly changes a precursor event in the successor structure
rather than fitting a scalar spatial representation. Across 25 independently
initialized strict channels, both deleting the activation and relocating its
post-transfer state prevented the same later articulation event. Exact address
placement is therefore causally important for this frozen outcome. The seven
memory-sensitive channels prevent a placement-only law: in those geometries,
activation and amplitude at the address were insufficient without some member
of the erased state bundle. The supported object is consequently a
coherence-conditioned transition geometry with heterogeneous marked-state
dependence, not a radial field or particle-like source.

The rebridge veto gives a complementary boundary. It usually prolonged the
topological opening, but never produced the preregistered 128-offset durable
separation. The result supports causal control over a successor-graph
transition; it does not establish durable morphogenesis or a physical
field-level correspondence.

Finding 075 then tested whether a frozen rooted local signature explained
susceptibility to that transition. The paired outcomes contained a suggestive
directional contrast, but the confirmatory probability gate narrowly failed and
the independent natural-closeness gate failed strongly. The registered result
is therefore negative. It rules out this bounded signature as a sufficient
account under the frozen swap policy; it does not rule out relational topology.
The next representation should preserve marked successor-event structure,
larger relational context, or temporal order rather than treating the observed
contrast as permission to tune the failed signature.

Finding 078 shows that structural placement and inherited marks can be
separated operationally and can interact non-additively in a marked one-step
successor representation. It does not show that this representation predicts
later articulation. Finding 079 was designed to supply that independent test,
but its support pipeline failed before outcome measurement. The correct status
of state-conditioned repair affordance is therefore untested, not negative.
This distinction preserves the positive causal placement result from Finding
074 while preventing developmental instrument behavior from becoming a
confirmation claim.

\subsection{Candidate bridge to Gravity from Entropy}
\label{sec:gfe-bridge}

The exact identity in Equation~\ref{eq:uniform-relative-entropy} supplies a
nonverbal bridge from the current support ratio to relative entropy. It does not
yet connect Iota dynamics to gravity. In GfE, the local Lagrangian is

\begin{equation}
\mathcal{L}_{\mathrm{GfE}}(x)
=-\operatorname{Tr}_{F}\ln\boldsymbol{\Theta}_x,
\qquad
\boldsymbol{\Theta}_x
=\widetilde{\mathbf{G}}_x\widetilde{g}_x^{-1},
\label{eq:gfe-lagrangian}
\end{equation}

where $\widetilde{g}$ is the physical higher-order metric,
$\widetilde{\mathbf{G}}$ is the metric induced by matter and curvature, and
$\boldsymbol{\Theta}$ is required to be positive
\cite{bianconi2025gravity,bianconi2026thermodynamics}. These metrics are not
normalized probability distributions. In the finite flattened representation,
Equation~\ref{eq:gfe-lagrangian} is a log-determinant contrast,
$\mathcal{L}_{\mathrm{GfE}}=-\ln\det_F\boldsymbol{\Theta}$.

A concrete bridge theorem would require a positive future-law operator
$\mathbf{Q}_x$ derived from the Iota transition structure such that

\begin{equation}
\begin{aligned}
\boldsymbol{\Theta}_x
&\stackrel{?}{=}
(\mathbf{Q}^{\mathrm{prior}}_x)^{-1/2}
\mathbf{Q}^{\mathrm{post}}_x
(\mathbf{Q}^{\mathrm{prior}}_x)^{-1/2}
\sim
\mathbf{Q}^{\mathrm{post}}_x
(\mathbf{Q}^{\mathrm{prior}}_x)^{-1},\\
\det_F\mathbf{Q}_x&=a\,\mu(\mathcal{R}_x),
\qquad a>0\ \text{common to prior and post}.
\end{aligned}
\label{eq:iota-gfe-operator-bridge}
\end{equation}

Here $\sim$ denotes similarity, so the first representative is positive while
retaining the determinant ratio of the original product. For the present finite
Iota identity, $\mu$ is counting measure. A non-counting extension would define
a new $\Delta I_x^{(\mu)}$ rather than silently replace the cardinality-based
$\Delta I_x$. Under these algebraic conditions,

\begin{equation}
\begin{aligned}
\mathcal{L}_{\mathrm{GfE}}(x)
&=\ln\frac{\mu(\mathcal{R}^{\mathrm{prior}}_x)}
{\mu(\mathcal{R}^{\mathrm{post}}_x)}\\
&=\ln\frac{|\mathcal{R}^{\mathrm{prior}}_x|}
{|\mathcal{R}^{\mathrm{post}}_x|}
=\Delta I_x.
\end{aligned}
\label{eq:iota-gfe-scalar-bridge}
\end{equation}

The question mark is substantive: a physical bridge must still derive the
operators, metric, covariance, variation, and continuum limit. Matching
determinants after the fact would only reproduce a scalar.

GfE thermodynamics further writes a local mode contribution as
$\delta s_k/\delta v=\ln\mathcal{G}_k$, so that
$\mathcal{G}_k=\exp(\delta s_k/\delta v)$ and its intrinsic geometric
temperature is $\theta_k=\mathcal{G}_k-1$
\cite{bianconi2026thermodynamics}. If an Iota mode-resolved contribution
$\chi_{x,k}$ and a physical mode dictionary were independently derived, the
conditional sectorwise crosswalk would be

\begin{equation}
\begin{aligned}
\chi_{x,k}&\widehat{=}c_k:=\frac{\delta s_k}{\delta v},\\
\mathcal{G}_k&=e^{\chi_{x,k}},\qquad
\theta_k=e^{\chi_{x,k}}-1,\qquad
\frac{\delta\epsilon_k}{\delta v}
=e^{\chi_{x,k}}-1-\chi_{x,k}.
\end{aligned}
\label{eq:gfe-conditional-crosswalk}
\end{equation}

Here $k\in\{0,\ldots,4\}$ labels the five GfE sectors and $z_k$ is the
corresponding degeneracy, with $z_0=z_1=1$ and $z_2=z_3=z_4=3$. In the
reported weak-curvature Friedmann expansion, $u_k=\omega_kH^2$ with
$\omega_k=\omega'_k\ell_{\mathrm P}^2$, and
$\bar\omega_{[1]}:=\sum_{k=0}^{4}z_k\omega_k$.

This cannot be imposed generally: individual GfE $c_k$ can be signed, whereas
Equation~\ref{eq:future-law-divergence} is nonnegative. Nor does an aggregate
$\chi_x\widehat{=}\sum_k z_kc_k$ determine a single temperature or energy,
because $\delta\epsilon/\delta v=\sum_kz_k(e^{c_k}-1-c_k)$ depends on the full
mode vector. The reference, normalization, mode dictionary, aggregation, and
sign structure must be derived rather than selected for agreement.

For weak-curvature, late-time Friedmann backgrounds, modes with
$\omega_k\neq0$ obey $c_k\simeq\omega_kH^2$ and
$\delta\epsilon_k/\delta v\simeq\tfrac12\omega_k^2H^4$. For non-de Sitter
solutions with $H\sim t^{-1}$ these scale as $t^{-2}$ and $t^{-4}$.
After summation, the leading entropy coefficient $\bar\omega_{[1]}$ can vanish,
including in the reported Milne and $w=1$ cases, so an aggregate $H^2$ scaling
requires $\bar\omega_{[1]}\neq0$. Under the conditional mode map, an Iota
quantity would inherit a scaling only in the corresponding nonvanishing sector.
That is not yet a novel Iota prediction: it is inherited from the Friedmann
approximation. Novel evidence would require an Iota-derived
correction, transition, or invariant tied to measured succession variables
such as opportunity lifetime or ordered cut depth, followed by a prediction
that defeats ordinary branching, percolation, complete-state dynamics, and
competing gravity models.

The bridge is therefore a precise research program. Construct
$\mathbf{Q}_x$ from policy-indexed future laws; prove positivity, covariance,
and robustness to scheduler, encoding, topology, and coarse-graining; recover
the GfE metric operator and action at the functional-variation level; establish
a controlled continuum limit; and only then derive an astrophysical or
cosmological discriminator. The current simulator results supply a causal
testbed for the first step, not evidence that the GfE action or gravity has
emerged.
